\setcounter{section}{4}

  \zwischenueberschrift{Pemetaan Weingarten}

Misalkan

\mavergleichskette
{\vergleichskette
{ Y
}
{ \subseteq }{ \R^n
}
{  }{
}
{    }{
}
{   }{
}
}
{}{}{}
adalah suatu
\definitionsverweis {hipermuka terdiferensialkan}{}{.}
Untuk suatu titik

\mavergleichskette
{\vergleichskette
{ P
}
{ \in }{ Y
}
{  }{
}
{    }{
}
{   }{
}
}
{}{}{}
kita mendefinisikan suatu pemetaan linear
\maabbdisp {L_P} { T_PY } { T_PY
} {}
pada ruang tangen $Y$ di P. Misalkan $N$ adalah suatu medan normal satuan terdiferensialkan pada $Y$ yang didefinisikan pada suatu lingkungan terbuka di sekitar $Y$. Gagasan pokoknya adalah merepresentasikan suatu vektor tangen

\mavergleichskette
{\vergleichskette
{ v
}
{ \in }{ T_PY
}
{  }{
}
{    }{
}
{   }{
}
}
{}{}{}
dengan suatu kurva terdiferensialkan
\maabbdisp {\gamma} { I } { Y
} {}
sehingga

\mavergleichskettedisp
{\vergleichskette
{ \gamma'(0)
}
{ =} { v
}
{ } {
}
{ } {
}
{ } {
}
}
{}{}{.}
Medan normal satuan itu kemudian mendefinisikan pemetaan
\maabbdisp {N \circ \gamma} { I } { \R^n
} {}
sepanjang $I$. Perubahan infinitesimal medan normal satuan sepanjang $\gamma$ diukur oleh limit

\mathdisp {\operatorname{lim}_{   t  \rightarrow  0   } \,  { \frac{   N(\gamma(t))-N(P)  }{  t  } }} {  }

jika limit ini ada. Menurut
Lemma 43.4 (Analysis (Osnabrück 2021-2023)),
limit ini sama dengan turunan berarah $\left(  D_{ v }  N   \right) \left(   P   \right)$ dan, pada gilirannya, sama dengan
\mathl{\left(  D N   \right)_{ P } \left(   v   \right)}{,} yaitu
\definitionsverweis {diferensial total}{}{}
yang dievaluasi pada vektor $v$. Akan ditunjukkan bahwa pemetaan
\maabbeledisp {} { T_PY } { \R^n
} {v } { \left(  D_{v}  N   \right) \left(   P   \right)
} {,}
ini mengambil nilai di $T_PY$.

\inputdefinition
{}
{

Misalkan

\mavergleichskette
{\vergleichskette
{ W
}
{ \subseteq }{ \R^n
}
{  }{
}
{    }{
}
{   }{
}
}
{}{}{}
\definitionsverweis {terbuka}{}{,}
\maabb {h} { W } { \R
} {}
suatu
\definitionsverweis {fungsi yang dua kali terdiferensialkan secara kontinu}{}{}
dan

\mavergleichskette
{\vergleichskette
{ Y
}
{ = }{ h^{-1}(c)
}
{  }{
}
{    }{
}
{   }{
}
}
{}{}{}
adalah
\definitionsverweis {serat}{}{}
di atas

\mavergleichskette
{\vergleichskette
{ c
}
{ \in }{ \R
}
{  }{
}
{    }{
}
{   }{
}
}
{}{}{,}
dengan $h$
\definitionsverweis {reguler}{}{}
pada setiap titik di $Y$. Misalkan $N$ adalah suatu
\definitionsverweis {medan normal satuan}{}{}
dan misalkan

\mavergleichskette
{\vergleichskette
{ P
}
{ \in }{ Y
}
{  }{
}
{    }{
}
{   }{
}
}
{}{}{.}
Pemetaan
\maabbeledisp {L_P} {T_P Y} { T_PY
} { v } { - \left(  D_{v}  N   \right) \left(   P  \right)
} {,}
disebut
\definitionswort {pemetaan Weingarten}{}
pada $T_PY$.

}

Perhatikan bahwa pemetaan Weingarten bergantung pada medan normal satuan, meskipun hal ini tidak selalu dinyatakan secara eksplisit. Jika $Y$ diberikan sebagai serat dari $h$, biasanya digunakan medan gradien ternormalisasi yang bersesuaian dengan $h$.


\inputfaktbeweis
{Hyperfläche/Weingartenabbildung/Tangentialraum/Fakt}
{Lemma}
{}
{

\faktsituation {Misalkan

\mavergleichskette
{\vergleichskette
{ W
}
{ \subseteq }{ \R^n
}
{  }{
}
{    }{
}
{   }{
}
}
{}{}{}
\definitionsverweis {terbuka}{}{,}
\maabb {h} { W } { \R
} {}
suatu
\definitionsverweis {fungsi yang dua kali terdiferensialkan secara kontinu}{}{}
dan

\mavergleichskette
{\vergleichskette
{ Y
}
{ = }{ h^{-1}(c)
}
{  }{
}
{    }{
}
{   }{
}
}
{}{}{}
adalah
\definitionsverweis {serat}{}{}
di atas

\mavergleichskette
{\vergleichskette
{ c
}
{ \in }{ \R
}
{  }{
}
{    }{
}
{   }{
}
}
{}{}{,}
dengan $h$
\definitionsverweis {reguler}{}{}
pada setiap titik di $Y$. Misalkan

\mavergleichskette
{\vergleichskette
{ P
}
{ \in }{ Y
}
{  }{
}
{    }{
}
{   }{
}
}
{}{}{.}}
\faktfolgerung {Maka
\definitionsverweis {pemetaan Weingarten}{}{}
$L_P$ adalah suatu endomorfisme linear dari
\definitionsverweis {ruang tangen}{}{}
$T_PY$.}
Catatan edisi: sumber memakai huruf kecil p pada ruang tangen ini, padahal titik yang dikuantifikasi adalah P.
\faktzusatz {}
\faktzusatz {}

}
{

Medan

\mavergleichskette
{\vergleichskette
{ N(P)
}
{ = }{ { \frac{
\operatorname{Grad} \,  h  (  P )  }{  \Vert {
\operatorname{Grad} \,  h  (  P ) } \Vert  } }
}
{  }{
}
{    }{
}
{   }{
}
}
{}{}{}
adalah suatu medan vektor yang terdiferensialkan secara kontinu pada suatu lingkungan terbuka

\mavergleichskette
{\vergleichskette
{ Y
}
{ \subseteq }{ U
}
{  }{
}
{    }{
}
{   }{
}
}
{}{}{}
tempat medan tersebut didefinisikan. Oleh karena itu, menurut
Proposition 46.1 (Analysis (Osnabrück 2021-2023)),

\mavergleichskettedisp
{\vergleichskette
{ \left(  D_{v}  N   \right) \left(   P  \right)
}
{ =} { \left(  D N   \right)_{ P } \left(  v  \right)
}
{ } {
}
{ } {
}
{ } {
}
}
{}{}{}
bersifat linear terhadap arah $v$. Dari syarat normal satuan diperoleh

\mavergleichskette
{\vergleichskette
{ \left\langle  N(P)  ,  N(P)   \right\rangle
}
{ = }{ 1
}
{  }{
}
{    }{
}
{   }{
}
}
{}{}{}
untuk setiap

\mavergleichskette
{\vergleichskette
{ P
}
{ \in }{ U
}
{  }{
}
{    }{
}
{   }{
}
}
{}{}{}
dan, dengan menggunakan
Soal 43.14 (Analysis (Osnabrück 2021-2023)), diperoleh

\mavergleichskettedisp
{\vergleichskette
{ 0
}
{ =} { \left(  D_{v}  \left\langle  N  ,  N  \right\rangle   \right) \left(   P  \right)
}
{ =} { 2 \left\langle  N(P)  ,  \left(  D_{v}  N   \right) \left(   P  \right)   \right\rangle
}
{ } {
}
{ } {
}
}
{}{}{.}
Jadi,
\mathl{\left(  D_{v}  N   \right) \left(   P  \right)}{} tegak lurus terhadap $N(P)$ dan, untuk

\mavergleichskette
{\vergleichskette
{ P
}
{ \in }{ Y
}
{  }{
}
{    }{
}
{   }{
}
}
{}{}{}
vektor tersebut termasuk dalam ruang tangen $T_PY$.

}

Dengan demikian, pemetaan Weingarten $L_P$ adalah negatif dari pembatasan diferensial total medan normal satuan pada ruang tangen $T_PY$, dengan kodomain ruang tangen $T_PY$. Jika diferensial total dihitung melalui
\definitionsverweis {matriks Jacobi}{}{}
dari $-N$, kita harus menentukan cara matriks itu bekerja pada suatu basis ruang tangen untuk memperoleh representasi matriks dari pemetaan Weingarten.



\inputfaktbeweis
{Ebene Kurve/Weingartenabbildung/Krümmung/Fakt}
{Lemma}
{}
{

\faktsituation {Misalkan

\mavergleichskette
{\vergleichskette
{ W
}
{ \subseteq }{ \R^2
}
{  }{
}
{    }{
}
{   }{
}
}
{}{}{}
\definitionsverweis {terbuka}{}{,}
\maabb {h} { W } { \R
} {}
suatu
\definitionsverweis {fungsi yang dua kali terdiferensialkan secara kontinu}{}{}
dan

\mavergleichskette
{\vergleichskette
{ C
}
{ = }{ h^{-1}(c)
}
{  }{
}
{    }{
}
{   }{
}
}
{}{}{}
adalah serat di atas

\mavergleichskette
{\vergleichskette
{ c
}
{ \in }{ \R
}
{  }{
}
{    }{
}
{   }{
}
}
{}{}{,}
dengan $h$
\definitionsverweis {reguler}{}{}
pada setiap titik di $C$, dan misalkan

\mavergleichskette
{\vergleichskette
{ P
}
{ \in }{ C
}
{  }{
}
{    }{
}
{   }{
}
}
{}{}{.}}
\faktfolgerung {Maka
\definitionsverweis {pemetaan Weingarten}{}{}
di $P$ adalah perkalian dengan
\definitionsverweis {kelengkungan}{}{}
$\kappa(P)$ dari $C$ di $P$.}
\faktzusatz {}
\faktzusatz {}

}
{

Pernyataan ini merupakan perumusan ulang dari
Lemma 3.11.

}

Pernyataan di atas tidak secara eksplisit merujuk pada orientasi; orientasi yang dimaksud harus turut diperhitungkan.

\inputbeispiel{}
{

Misalkan

\mavergleichskettedisp
{\vergleichskette
{ h(x,y,z)
}
{ =} { x^2+y^2+z^2
}
{ } {
}
{ } {
}
{ } {
}
}
{}{}{}
dan perhatikan serat $Y$ dari $h$ di atas $r^2$, yakni permukaan bola berjari-jari

\mavergleichskette
{\vergleichskette
{ r
}
{ > }{ 0
}
{  }{
}
{    }{
}
{   }{
}
}
{}{}{}
yang berpusat di titik asal. Misalkan

\mavergleichskette
{\vergleichskette
{ P
}
{ = }{ \left(  x_0   , \,   y_0  , \,  z_0                 \right)
}
{ \in }{ Y
}
{    }{
}
{   }{
}
}
{}{}{.}
Dengan suatu isometri, titik ini dapat dipetakan ke

\mavergleichskette
{\vergleichskette
{ Q
}
{ = }{ \left(  r   , \,   0  , \,   0                 \right)
}
{  }{
}
{    }{
}
{   }{
}
}
{}{}{}
tanpa mengubah
\definitionsverweis {pemetaan Weingarten}{}{}.
Suatu basis ruang tangennya ialah
\mathkor {} {\begin{pmatrix}  0  \\ 1 \\  0   \end{pmatrix}} {dan} {\begin{pmatrix}  0  \\ 0\\  1   \end{pmatrix}} {.}
Medan normal satuan $N$ yang mengarah ke dalam adalah
\mathl{- { \frac{   1  }{  2r } } \begin{pmatrix}  2x \\ 2y\\  2z   \end{pmatrix}}{} dan karena itu, untuk

\mavergleichskettedisp
{\vergleichskette
{ v
}
{ =} { \begin{pmatrix}  0  \\ b \\ c   \end{pmatrix}
}
{ } {
}
{ } {
}
{ } {
}
}
{}{}{}

\mavergleichskettedisp
{\vergleichskette
{ \left(  D_{v}  N   \right) \left(   Q   \right)
}
{ =} { b { \frac{   \partial   N   }{  \partial   y  } } (Q) +  c { \frac{   \partial   N   }{  \partial  z  } }  (Q)
}
{ =} { -b { \frac{   1  }{  2r } } \begin{pmatrix}  0  \\ 2 \\  0   \end{pmatrix} - c { \frac{   1  }{  2r } } \begin{pmatrix}  0  \\ 0\\  2   \end{pmatrix}
}
{ =} { - { \frac{   1  }{ r } } \begin{pmatrix}  0  \\ b \\ c   \end{pmatrix}
}
{ } {
}
}
{}{}{.}
Jadi, pemetaan Weingarten adalah perkalian dengan kebalikan ${ \frac{   1  }{ r } }$ dari jari-jari.

}

\inputbeispiel{}
{

Kita melanjutkan
Contoh 2.5.
Matriks Jacobi dari medan normal satuan

\mavergleichskettedisphandlinks
{\vergleichskettedisphandlinks
{ N( \begin{pmatrix}  x  \\ y \\ z   \end{pmatrix} )
}
{ =} { \left(  x^2+y^2+z^2  \right)^{ - { \frac{   1  }{  2 } } } \begin{pmatrix}  x  \\ y \\  -z   \end{pmatrix}
}
{ } {
}
{ } {
}
{ } {
}
}
{}{}{}
adalah

\mathdisp {\left(  x^2+y^2+z^2  \right)^{- { \frac{   3  }{  2 } } } \begin{pmatrix}  y^2+z^2  &   -xy &  -xz \\  -xy &  x^2+z^2 &  -yz \\ xz &  yz &  -x^2-y^2  \end{pmatrix}} { . }

\definitionsverweis {Pemetaan Weingarten}{}{}
adalah negatif dari pembatasan pemetaan ini pada ruang tangen permukaan tersebut; menurut
Lemma 4.2,
hasilnya kembali berada di ruang tangen. Kita mengasumsikan

\mavergleichskette
{\vergleichskette
{ x
}
{ \neq }{ 0
}
{  }{
}
{    }{
}
{   }{
}
}
{}{}{}
dan menggunakan basis
\mathkor {} {\begin{pmatrix}  y  \\ -x\\  0   \end{pmatrix}} {dan} {\begin{pmatrix}  z  \\ 0 \\  x   \end{pmatrix}} {}
dari ruang tangen. Kita mempunyai

\mavergleichskettealignhandlinks
{\vergleichskettealignhandlinks
{ \begin{pmatrix}  y^2+z^2  &   -xy &  -xz \\  -xy &  x^2+z^2 &  -yz \\ xz &  yz &  -x^2-y^2  \end{pmatrix} \begin{pmatrix}  y  \\ -x\\  0   \end{pmatrix}
}
{ =} { \begin{pmatrix}  y \left(  y^2+z^2  \right) +x^2y \\ -xy^2 -x \left(  x^2+z^2  \right) \\  0   \end{pmatrix}
}
{ =} { \left(  x^2+y^2+z^2  \right) \begin{pmatrix}  y  \\ -x\\  0   \end{pmatrix}
}
{ } {
}
{ } {
}
}
{}
{}{,}
sehingga langsung diperoleh vektor eigen $\begin{pmatrix}  y  \\ -x\\  0   \end{pmatrix}$ dengan nilai eigen
\mathl{- \left(  x^2+y^2+z^2  \right)^{ - { \frac{   1  }{  2 } } }}{}.
Selain itu,


\mavergleichskettealigndrucklinks
{\vergleichskettealigndrucklinks
{ \begin{pmatrix}  y^2+z^2  &   -xy &  -xz \\  -xy &  x^2+z^2 &  -yz \\ xz &  yz &  -x^2-y^2  \end{pmatrix} \begin{pmatrix}  z  \\ 0 \\  x   \end{pmatrix}
}
{ =} { \begin{pmatrix} z \left(  y^2+z^2  \right) -x^2z \\ -2xyz\\  xz^2-x \left(  x^2+y^2  \right)   \end{pmatrix}
}
{ =} { 2yz \begin{pmatrix}  y  \\ -x\\  0   \end{pmatrix} + \left(  z^2-x^2-y^2  \right) \begin{pmatrix}  z  \\ 0 \\  x   \end{pmatrix}
}
{ =} { 2yz \begin{pmatrix}  y  \\ -x\\  0   \end{pmatrix}- \begin{pmatrix}  z  \\ 0 \\  x   \end{pmatrix}
}
{ } {
}
}
{}{}{.}

Dengan demikian, terhadap basis ini, pemetaan Weingarten direpresentasikan oleh matriks

\mathdisp {\begin{pmatrix}  - \left(  x^2+y^2+z^2  \right)^{ - { \frac{   1  }{  2 } } }   &   - 2yz \left(  x^2+y^2+z^2  \right)^{ - { \frac{   3  }{  2 } } }    \\  0  &  \left(  x^2+y^2+z^2  \right)^{ - { \frac{   3  }{  2 } } }   \end{pmatrix}} {  }

Untuk memperoleh vektor eigen kedua dalam ruang tangen,
\zusatzklammer {dengan memperhatikan
Korolari 4.8} {} {}
cara termudah ialah menentukan suatu vektor yang tegak lurus terhadap vektor eigen pertama dan vektor normal. Hasilnya adalah vektor
\mathl{\begin{pmatrix}  xz \\ yz\\  x^2+y^2   \end{pmatrix}}{,} dan memang,
\zusatzklammer {tanpa faktor skalar di depan} {} {}


\mavergleichskettealigndrucklinks
{\vergleichskettealigndrucklinks
{ \begin{pmatrix}  y^2+z^2  &   -xy &  -xz \\  -xy &  x^2+z^2 &  -yz \\ xz &  yz &  -x^2-y^2  \end{pmatrix} \begin{pmatrix}  xz \\ yz\\  x^2+y^2   \end{pmatrix}
}
{ =} { \begin{pmatrix}  xz \left(  y^2+z^2  \right) -xy^2z -xz \left(  x^2+y^2  \right)  \\ - x^2yz +yz \left(  x^2+z^2  \right) -yz \left(  x^2+y^2  \right) \\  x^2z^2 +y^2z^2 - \left(  x^2+y^2  \right)^2   \end{pmatrix}
}
{ =} { \left(  - x^2-y^2+z^2  \right) \begin{pmatrix}  xz \\ yz\\  x^2+y^2   \end{pmatrix}
}
{ =} { - \begin{pmatrix}  xz \\ yz\\  x^2+y^2   \end{pmatrix}
}
{ } {
}
}
{}
{}{.}

Jadi, nilai eigennya adalah
\mathl{\left(  x^2+y^2+z^2  \right)^{- { \frac{   3  }{  2 } } }}{}
\zusatzklammer {hal ini juga dapat langsung dibaca dari matriks segitiga yang merepresentasikannya} {} {.}

}




\inputfaktbeweis
{Hyperfläche/Weingartenabbildung/Zweite Ableitung/Skalarprodukt/Fakt}
{Lemma}
{}
{

\faktsituation {Misalkan

\mavergleichskette
{\vergleichskette
{ W
}
{ \subseteq }{ \R^n
}
{  }{
}
{    }{
}
{   }{
}
}
{}{}{}
\definitionsverweis {terbuka}{}{,}
\maabb {h} { W } { \R
} {}
suatu
\definitionsverweis {fungsi yang dua kali terdiferensialkan secara kontinu}{}{}
dan

\mavergleichskette
{\vergleichskette
{ Y
}
{ = }{ h^{-1}(c)
}
{  }{
}
{    }{
}
{   }{
}
}
{}{}{}
adalah
\definitionsverweis {serat}{}{}
di atas

\mavergleichskette
{\vergleichskette
{ c
}
{ \in }{ \R
}
{  }{
}
{    }{
}
{   }{
}
}
{}{}{,}
dengan $h$
\definitionsverweis {reguler}{}{}
pada setiap titik di $Y$. Misalkan $N$ adalah suatu
\definitionsverweis {medan normal satuan}{}{}
dan misalkan

\mavergleichskette
{\vergleichskette
{ P
}
{ \in }{ Y
}
{  }{
}
{    }{
}
{   }{
}
}
{}{}{.}
Misalkan $\gamma$ adalah suatu realisasi yang terdiferensialkan dua kali pada $Y$ dari vektor tangen

\mavergleichskette
{\vergleichskette
{ v
}
{ \in }{ T_PY
}
{  }{
}
{    }{
}
{   }{
}
}
{}{}{.}}
\faktfolgerung {Maka

\mavergleichskettedisp
{\vergleichskette
{ \left\langle  \gamma^{\prime \prime} (0)  ,  N(P)  \right\rangle
}
{ =} { \left\langle  L_P(v) , v  \right\rangle
}
{ } {
}
{ } {
}
{ } {
}
}
{}{}{.}}
\faktzusatz {}
\faktzusatz {}

}
{

Misalkan
\maabbeledisp {\gamma} { I } { Y
} { t } { \gamma(t)
} {,}
terdiferensialkan dua kali, dengan

\mavergleichskette
{\vergleichskette
{ \gamma(0)
}
{ = }{ P
}
{  }{
}
{    }{
}
{   }{
}
}
{}{}{}
dan

\mavergleichskette
{\vergleichskette
{ \gamma'(0)
}
{ = }{ v
}
{ \in }{ T_PY
}
{    }{
}
{   }{
}
}
{}{}{.}
Kita mempunyai

\mavergleichskettedisp
{\vergleichskette
{ \left\langle  \gamma'(t)  ,  N(\gamma(t))   \right\rangle
}
{ =} { 0
}
{ } {
}
{ } {
}
{ } {
}
}
{}{}{}
untuk semua

\mavergleichskette
{\vergleichskette
{ t
}
{ \in }{ I
}
{  }{
}
{    }{
}
{   }{
}
}
{}{}{}
sehingga, dengan
Soal 43.14 (Analysis (Osnabrück 2021-2023)), diperoleh

\mavergleichskettealign
{\vergleichskettealign
{ 0
}
{ =} { \left\langle  \gamma'(t)  ,  N(\gamma(t))   \right\rangle' (0)
}
{ =} { \left\langle  \gamma^{\prime \prime}(0)  ,  N(\gamma(0))  \right\rangle + \left\langle  \gamma'(0)  ,  \left(  D_{ \gamma'(0) }  N   \right) \left(   \gamma(0)   \right)   \right\rangle
}
{ =} { \left\langle  \gamma^{\prime \prime}(0) , N(P)  \right\rangle - \left\langle  L_P(v)  ,  v   \right\rangle
}
{ } {
}
}
{}
{}{.}

}

Dalam pernyataan di atas tidak diperlukan penetapan tambahan mengenai orientasi, sebab medan normal satuan muncul pada kedua ruas, dan pada ruas kanan muncul melalui pemetaan Weingarten. Perhatikan pula bahwa $\gamma$ tidak harus memiliki kecepatan konstan. Dengan demikian, terdapat banyak realisasi dan percepatan $\gamma^{\prime \prime} (0)$ tidak ditentukan secara tunggal. Namun, nilai
\mathl{\left\langle  \gamma^{\prime \prime}(0)  ,  N(P)  \right\rangle}{,} yang menggambarkan komponen normal percepatan, tidak berubah karena $\gamma$ bergerak pada hipermuka tersebut.
Catatan edisi: sumber kehilangan argumen N(P) dalam hasil kali dalam pada kalimat sebelumnya; identitas yang baru saja dibuktikan menentukannya secara unik.



\inputfaktbeweis
{Hyperfläche/Weingartenabbildung/Selbstadjungiert/Fakt}
{Satz}
{}
{

\faktsituation {Misalkan

\mavergleichskette
{\vergleichskette
{ W
}
{ \subseteq }{ \R^n
}
{  }{
}
{    }{
}
{   }{
}
}
{}{}{}
\definitionsverweis {terbuka}{}{,}
\maabb {h} { W } { \R
} {}
suatu
\definitionsverweis {fungsi yang dua kali terdiferensialkan secara kontinu}{}{}
dan

\mavergleichskette
{\vergleichskette
{ Y
}
{ = }{ h^{-1}(c)
}
{  }{
}
{    }{
}
{   }{
}
}
{}{}{}
adalah serat di atas

\mavergleichskette
{\vergleichskette
{ c
}
{ \in }{ \R
}
{  }{
}
{    }{
}
{   }{
}
}
{}{}{,}
dengan $h$
\definitionsverweis {reguler}{}{}
pada setiap titik di $Y$.}
\faktfolgerung {Maka
\definitionsverweis {pemetaan Weingarten}{}{}
$L_P$ pada setiap titik

\mavergleichskette
{\vergleichskette
{ P
}
{ \in }{ Y
}
{  }{
}
{    }{
}
{   }{
}
}
{}{}{}
bersifat
\definitionsverweis {adjoin-diri (self-adjoint)}{}{.}}
\faktzusatz {}
\faktzusatz {}

}
{

Untuk vektor

\mavergleichskette
{\vergleichskette
{ v,w
}
{ \in }{ T_PY
}
{  }{
}
{    }{
}
{   }{
}
}
{}{}{}
kita harus menunjukkan bahwa

\mavergleichskettedisp
{\vergleichskette
{ \left\langle  v  ,  L_P(w)  \right\rangle
}
{ =} { \left\langle L_P(v) ,  w  \right\rangle
}
{ } {
}
{ } {
}
{ } {
}
}
{}{}{}
tersebut. Dengan

\mavergleichskette
{\vergleichskette
{ N(P)
}
{ = }{ { \frac{
\operatorname{Grad} \,  h  (  P )  }{  \Vert {
\operatorname{Grad} \,  h  (  P ) } \Vert  } }
}
{  }{
}
{    }{
}
{   }{
}
}
{}{}{}
menurut
Lemma 45.10 (Analysis (Osnabrück 2021-2023)), kita mempunyai

\mavergleichskettealignhandlinks
{\vergleichskettealignhandlinks
{ \left\langle  v  ,  L_P(w)  \right\rangle
}
{ =} { -\left\langle  v  ,  \left(  D_{w}  N   \right) \left(   P   \right)   \right\rangle
}
{ =} { -\left\langle  v  ,  \left(  D_{w}  { \frac{
\operatorname{Grad} \,  h  }{  \Vert {
\operatorname{Grad} \,  h } \Vert  } }   \right) \left(   P   \right)     \right\rangle
}
{ =} { -\left\langle  v  ,  \left(  D_{w}  { \frac{   1  }{  \Vert {
\operatorname{Grad} \,  h } \Vert  } }   \right) \left(   P   \right) \cdot
\operatorname{Grad} \,  h  (  P ) + { \frac{   1  }{  \Vert {
\operatorname{Grad} \,  h  (  P ) } \Vert  } } \cdot  \left(  D_{w}
\operatorname{Grad} \,  h   \right) \left(   P   \right)   \right\rangle
}
{ =} { - { \frac{   1  }{  \Vert {
\operatorname{Grad} \,  h  (  P ) } \Vert  } } \left\langle  v  ,  \left(  D_{w}
\operatorname{Grad} \,  h   \right) \left(   P   \right)    \right\rangle
}
}
{}
{}{,}
karena suku pertama tegak lurus terhadap vektor tangen $v$. Dalam fungsi-fungsi koordinat,

\mavergleichskettealignhandlinks
{\vergleichskettealignhandlinks
{ \left(  D_{w}
\operatorname{Grad} \,  h   \right) \left(   P   \right)
}
{ =} { \left(  D_{w}  \left(  { \frac{   \partial   h   }{  \partial   x_1   } }   , \,   \ldots  , \,   { \frac{   \partial   h   }{  \partial   x_n   } }                 \right)    \right) \left(   P   \right)
}
{ =} { \sum_{j = 1}^n  w_j  { \frac{   \partial    }{  \partial   x_j   } }  \left(  { \frac{   \partial   h   }{  \partial   x_1   } }   , \,   \ldots  , \,   { \frac{   \partial   h   }{  \partial   x_n   } }                 \right) (P)
}
{ =} { \left(  \sum_{j = 1}^n  w_j  { \frac{   \partial    }{  \partial   x_j   } }  { \frac{   \partial   h   }{  \partial   x_1   } } (P)   , \,   \ldots  , \,   \sum_{j = 1}^n  w_j  { \frac{   \partial    }{  \partial   x_j   } }  { \frac{   \partial   h   }{  \partial   x_n   } } (P)                 \right)
}
{ } {
}
}
{}
{}{.}
Dengan demikian, ungkapan di atas sama dengan

\mavergleichskettealign
{\vergleichskettealign
{ \left\langle  v  ,  L_P(w)  \right\rangle
}
{ =} { - { \frac{   1  }{  \Vert {
\operatorname{Grad} \,  h  (  P ) } \Vert  } } \left\langle  v  ,  \left(  D_{w}
\operatorname{Grad} \,  h   \right) \left(   P   \right)   \right\rangle
}
{ =} { -  { \frac{   1  }{  \Vert {
\operatorname{Grad} \,  h  (  P ) } \Vert  } }  \sum_{i,j = 1}^n v_i  w_j  { \frac{   \partial    }{  \partial   x_j   } } { \frac{   \partial   h   }{  \partial   x_i   } }  (P)
}
{ } {
}
{ } {
}
}
{}
{}{.}
Menurut
Teorema 44.10 (Analysis (Osnabrück 2021-2023)),
urutan turunan parsial dapat dipertukarkan, sehingga peran
\mathkor {} {v} {dan} {w} {}
juga dapat dipertukarkan.

}



\inputfaktbeweis
{Hyperfläche/Weingartenabbildung/Selbstadjungiert/Eigenwerte/Fakt}
{Korollar}
{}
{

\faktsituation {Misalkan

\mavergleichskette
{\vergleichskette
{ W
}
{ \subseteq }{ \R^n
}
{  }{
}
{    }{
}
{   }{
}
}
{}{}{}
\definitionsverweis {terbuka}{}{,}
\maabb {h} { W } { \R
} {}
suatu
\definitionsverweis {fungsi yang dua kali terdiferensialkan secara kontinu}{}{}
dan

\mavergleichskette
{\vergleichskette
{ Y
}
{ = }{ h^{-1}(c)
}
{  }{
}
{    }{
}
{   }{
}
}
{}{}{}
adalah serat di atas

\mavergleichskette
{\vergleichskette
{ c
}
{ \in }{ \R
}
{  }{
}
{    }{
}
{   }{
}
}
{}{}{,}
dengan $h$
\definitionsverweis {reguler}{}{}
pada setiap titik di $Y$.}
\faktfolgerung {Maka
\definitionsverweis {pemetaan Weingarten}{}{}
$L_P$ pada setiap titik

\mavergleichskette
{\vergleichskette
{ P
}
{ \in }{ Y
}
{  }{
}
{    }{
}
{   }{
}
}
{}{}{}
\definitionsverweis {dapat didiagonalkan}{}{}
dengan
\definitionsverweis {nilai-nilai eigen}{}{}
real dan ruang-ruang eigen yang saling ortogonal.}
\faktzusatz {}
\faktzusatz {}

}
{

Hal ini mengikuti dari
Teorema 4.7
dan
Teorema 41.11 (Lineare Algebra (Osnabrück 2024-2025)).

}



\inputfaktbeweis
{Differenzierbare Funktion/Graph/Hyperfläche/Weingartenabbildung/Fakt}
{Lemma}
{}
{

\faktsituation {Misalkan

\mavergleichskette
{\vergleichskette
{ V
}
{ \subseteq }{ \R^{n-1}
}
{  }{
}
{    }{
}
{   }{
}
}
{}{}{}
suatu
\definitionsverweis {himpunan terbuka}{}{}
dan misalkan

\mavergleichskette
{\vergleichskette
{ Y
}
{ \subseteq }{ V \times \R
}
{  }{
}
{    }{
}
{   }{
}
}
{}{}{}
adalah
\definitionsverweis {grafik}{}{}
dari suatu fungsi yang
\definitionsverweis {dua kali terdiferensialkan secara kontinu}{}{}
\maabbdisp {f} { V } {\R
} {.}}
\faktfolgerung {Maka
\definitionsverweis {pemetaan Weingarten}{}{}
$L_P$ pada suatu titik

\mavergleichskette
{\vergleichskette
{ P
}
{ = }{ (x_1 , \ldots , x_{n-1}, f(x_1 , \ldots , x_{n-1}))
}
{ \in }{ Y
}
{    }{
}
{   }{
}
}
{}{}{}
\zusatzklammer {terhadap medan normal satuan yang mengarah ke atas} {} {}
dapat direpresentasikan sebagai berikut. Tuliskan \mathl{g=\operatorname{Grad} f}{} sebagai vektor kolom dan
\mathl{G=I_{n-1}+gg^{\mathsf T}}{,}
yakni matriks metrik dalam basis grafik. Matriks pemetaan Weingarten dalam basis tersebut adalah
\mathl{G^{-1}{ \frac{   1  }{  \sqrt{ 1+ \left(  { \frac{   \partial   f   }{  \partial   x_1   } }  \right)^2 + \cdots + \left(  { \frac{   \partial   f   }{  \partial   x_{n-1}   } }  \right)^2 }  } } \operatorname{Hess}_{   } \,  f}{}, dengan $\operatorname{Hess}_{   } \,  f$ menyatakan
\definitionsverweis {matriks Hesse}{}{}
dari $f$, apabila vektor-vektor dasar

\mavergleichskette
{\vergleichskette
{ v
}
{ \in }{ \R^{n-1}
}
{  }{
}
{    }{
}
{   }{
}
}
{}{}{}
diidentifikasi dengan vektor-vektor tangen dari $T_PY$ dalam pengertian
Contoh 1.2.}
\faktzusatz {}
\faktzusatz {}

}
{

Kita melanjutkan
Contoh 1.2.
Medan normal satuan yang mengarah ke atas diberikan oleh

\mavergleichskettedisphandlinks
{\vergleichskettedisphandlinks
{ N(x_1  , \ldots , x_{n-1}, f(x_1 , \ldots , x_{n-1}))
}
{ =} { { \frac{   1  }{  \sqrt{ 1+ \left(  { \frac{   \partial   f   }{  \partial   x_1   } }  \right)^2 + \cdots + \left(  { \frac{   \partial   f   }{  \partial   x_{n-1}   } }  \right)^2 }  } } \begin{pmatrix}  -{ \frac{   \partial   f   }{  \partial   x_{1}   } }  \\\vdots\\  -{ \frac{   \partial   f   }{  \partial   x_{n-1}   } }\\ 1   \end{pmatrix}
}
{ } {
}
{ } {
}
{ } {
}
}
{}{}{}
yang akan digunakan. Kita meninjau lintasan

\mavergleichskettedisp
{\vergleichskette
{ \gamma(t)
}
{ =} { \begin{pmatrix}  x_1+tv_1  \\ \vdots \\  x_{n-1} +tv_{n-1}\\f \begin{pmatrix}  x_1+tv_1  \\ \vdots \\  x_{n-1} +tv_{n-1}   \end{pmatrix}   \end{pmatrix}
}
{ } {
}
{ } {
}
{ } {
}
}
{}{}{}
pada grafik yang bersesuaian dengan vektor dasar $v$. Turunan keduanya adalah

\mavergleichskettedisp
{\vergleichskette
{ \gamma^{\prime \prime} (0)
}
{ =} { \begin{pmatrix}  0  \\ \vdots \\  0\\ \sum_{1 \leq i, j \leq n-1}  { \frac{   \partial    }{  \partial   x_i   } }  { \frac{   \partial   f   }{  \partial   x_j   } } v_iv_j   \end{pmatrix}
}
{ } {
}
{ } {
}
{ } {
}
}
{}{}{.}
Menurut
Lemma 4.6,

\mavergleichskettealigndrucklinks
{\vergleichskettealigndrucklinks
{ \left\langle  L_P( \gamma'(0) ) ,  \gamma'(0)  \right\rangle
}
{ =} { \left\langle  \gamma^{\prime \prime} (0) ,  N(x_1 , \ldots , x_{n-1}, f(x_1  , \ldots , x_{n-1}))   \right\rangle
}
{ =} { { \frac{   1  }{  \sqrt{ 1 + \left(  { \frac{   \partial   f   }{  \partial   x_1   } }  \right)^2 + \cdots + \left(  { \frac{   \partial   f  }{  \partial   x_{n-1}   } }  \right)^2 }  } } \left\langle  \begin{pmatrix}  0  \\\vdots\\  0\\ \sum_{ 1 \leq i, j \leq n-1 }  { \frac{   \partial    }{  \partial   x_i   } }  { \frac{   \partial   f   }{  \partial   x_j   } } v_iv_j   \end{pmatrix}  ,  \begin{pmatrix}  -{ \frac{   \partial   f   }{  \partial   x_{1}   } }  \\\vdots\\  -{ \frac{   \partial   f   }{  \partial   x_{n-1}   } }\\ 1   \end{pmatrix}   \right\rangle
}
{ =} { { \frac{   \sum_{ 1 \leq i, j \leq n-1 } { \frac{   \partial    }{  \partial   x_i   } }  { \frac{   \partial   f   }{  \partial   x_j   } } v_i v_j     }{  \sqrt{ 1+  \left(  { \frac{   \partial   f   }{  \partial   x_1   } }  \right)^2 + \cdots + \left(  { \frac{   \partial   f  }{  \partial   x_{n-1}   } }  \right)^2 }  } }
}
{ =} { { \frac{   1  }{  \sqrt{ 1+ \left(  { \frac{   \partial   f   }{  \partial   x_1   } }  \right)^2 + \cdots + \left(  { \frac{   \partial   f   }{  \partial   x_{n-1}   } }  \right)^2 }  } }  \left( v_1   , \,   \ldots  , \,   v_{n-1}                 \right)  \operatorname{Hess}_{   } \,  f  \begin{pmatrix} v_1  \\ \vdots \\  v_{n-1}   \end{pmatrix}
}
}
{}{}{.}
Artinya, menurut
Teorema 4.7,
\definitionsverweis {bentuk bilinear}{}{}
simetris

\mathl{\left\langle  L_P(-) ,  -  \right\rangle}{} pada ruang tangen $T_PY$ bersesuaian dengan bentuk bilinear yang diberikan oleh matriks Hesse yang diskalakan
\zusatzklammer {oleh faktor di depan} {} {}
pada $\R^{n-1}$, apabila vektor yang sama disubstitusikan pada kedua tempat.
Menurut
Lemma 38.10 (Lineare Algebra (Osnabrück 2024-2025)),
kedua bentuk bilinear tersebut berimpit untuk sembarang pasangan argumen. Akan tetapi, ini menentukan bentuk fundamental kedua, bukan langsung matriks operator dalam basis grafik. Jika \mathl{A}{} adalah matriks $L_P$ dalam basis grafik dan \mathl{G=I_{n-1}+gg^{\mathsf T}}{} adalah matriks metriknya, maka matriks bentuk bilinear di atas adalah \mathl{GA}{.} Oleh karena itu,

\mathdisp {GA={ \frac{   1  }{  \sqrt{1+\lVert g\rVert^2}  } }\operatorname{Hess}_{   }\, f} { , }

sehingga

\mathdisp {A=G^{-1}{ \frac{   1  }{  \sqrt{1+\lVert g\rVert^2}  } }\operatorname{Hess}_{   }\, f} { . }

Catatan edisi: sumber menghilangkan faktor invers metrik G dan dengan demikian menyamakan matriks bentuk fundamental kedua dengan matriks pemetaan Weingarten. Rumus dan simpulan di atas memperbaiki pembedaan tersebut.

}
